% ORA0 VA Compound — CCF-A conference paper (LaTeX)
% Compiled with pdflatex; figures to be filled from actual run logs.
\documentclass[10pt,twocolumn]{article}

\usepackage[margin=0.9in]{geometry}
\usepackage{amsmath,amssymb,amsfonts}
\usepackage{graphicx}
\usepackage{booktabs}
\usepackage{multirow}
\usepackage{hyperref}
\usepackage{microtype}
\usepackage[table]{xcolor}
\usepackage{caption}
\captionsetup{font=small,labelfont=bf}

\newcommand{\LFM}{\mathcal{L}_{\mathrm{FM}}}
\newcommand{\Lpair}{\mathcal{L}_{\mathrm{pair}}}
\newcommand{\todo}[1]{\textcolor{red}{[TODO: #1]}}

\title{Bidirectional Visual--Action Memory with Frozen Language Cache:
Structural Language Grounding for Lightweight VLA Policies}

\author{Anonymous Submission}

\begin{document}
\maketitle

\begin{abstract}
Vision-Language-Action (VLA) models inherit instruction-following ability from
large pretrained backbones, but policy fine-tuning erodes it through two
mechanisms: \emph{representation erosion}---action gradients rewriting
language-related parameters (instruction blindness)---and \emph{visual
shortcut learning}---the policy ignoring language because scene features
suffice. We show that a \emph{decoupled} architecture sidesteps the first
mechanism by construction: a frozen language backbone encodes the instruction
once into a static key/value cache, and a compact recurrent
``visual--action (VA) composite'' reads that cache through shared attention,
so action gradients can never reach language parameters. On LIBERO (3-scene,
12-task) and MetaWorld MT50, the frozen-cache policy retains sharp language
grounding (blanking language changes chunk error by +2381\%; wrong
instructions +108.5\%), while the same policy with an end-to-end fine-tuned
language adapter collapses the instruction embedding space (pairwise cosine
0.999 vs.\ 0.765) and loses language sensitivity (+1.5\%). A $2\times2$
control (frozen/LoRA language $\times$ frozen/unfrozen vision) isolates the
collapse to the trainable language adapter. The trainable part of our policy
is 43.5M parameters and deploys at 40.6~Hz with a two-term loss
(flow matching + paired-instruction contrast), no auxiliary objectives.
\todo{L\_m verdict, VLA-RL gains, MW multi-start, LIBERO-100}
\end{abstract}

\section{Introduction}
\label{sec:intro}

VLA models have become the dominant paradigm for robot manipulation: a
pretrained vision-language backbone conditions a continuous action head on the
visual observation and a natural-language instruction. Their promise is
\emph{language grounding}---executing the specific instruction the user gave,
rather than a scene-dependent default. Yet recent analyses show this ability
is fragile: fine-tuning a VLA on action data progressively degrades the
instruction-following behavior of its backbone
\cite{libero_cf,zhang2026flatnesspreservesinstructionfollowing}. Two mechanisms are implicated. First,
\emph{representation erosion}: action-loss gradients flow back through the
shared backbone and rewrite the parameters carrying language meaning,
collapsing the model's ability to distinguish instructions. Second,
\emph{visual shortcut learning}: because scene features are highly predictive
of the correct action in most datasets, the policy can minimize training loss
while attending to language almost not at all---it ``passes'' on the data
distribution but silently fails on counterfactual instructions. Both
mechanisms are \emph{training-coupled}: they arise from the interaction
between the action loss and a shared, trainable backbone.

In this paper we ask whether a \emph{decoupled} architecture can remove the
first mechanism by construction and expose the second through controlled
counterfactual evaluation. We study a compact policy, the
\emph{bidirectional visual--action (VA) composite}: a frozen language encoder
(Qwen-3.5-2B) runs once per instruction and caches its final-layer hidden
states as static key/value anchors; a frozen video encoder (V-JEPA 2.1) turns
a causal 4-frame window into visual tokens; and a small stack of shared
attention blocks---the VA composite---reads vision, a learned action stream,
a single-slot recurrent visual memory and the static language anchors, before
a flow-matching head regresses the action chunk. Because the language encoder
is frozen \emph{and} detached, action gradients cannot reach language
parameters by construction: representation erosion is structurally
impossible, not merely discouraged.

Our contributions:
\begin{enumerate}
\item \textbf{Architecture-level immunity to representation erosion.} A
$2\times2$ control (language frozen/LoRA-adapted $\times$ vision
frozen/unfrozen) on LIBERO shows: a fully frozen pipeline retains the exact
pretrained instruction embedding space (pairwise cosine of 12 instructions:
0.765, identical to the untrained encoder), while end-to-end fine-tuning with
a LoRA adapter collapses that space (0.999) and destroys behavioral language
sensitivity (blanking the language stream changes chunk error by +1.5\%
vs.\ +2381\% for the frozen policy).
\item \textbf{Causal language-grounding evidence under same-scene
counterfactuals.} On LIBERO's same-scene different-instruction regime---the
regime that exposes visual shortcuts---language is a necessary input
(+2381\% blank sensitivity, 3-scene; +13751\% 1-scene). On MetaWorld MT50,
wrong instructions degrade action error by +108.5\%, task-id substitution by
+81.5\%.
\item \textbf{Lightweight and simple.} 43.5M trainable parameters, 40.6~Hz
deployment on a single consumer GPU, two-term loss
(flow matching + paired-instruction contrast) with no auxiliary objectives.
\item \textbf{Honest scaling report.} We trace our current closed-loop gap on
MetaWorld to training-data coverage (samples cover only the first 0.33~s of
each episode), verified by in-distribution tasks succeeding at ceiling.
\todo{multi-start results}
\end{enumerate}

We emphasize the scope of our claims: we do not argue that end-to-end VLA
fine-tuning universally destroys semantic representations---only that within
our tested configurations (V-JEPA 2.1 ViT-B, Qwen 2B, LIBERO), the frozen-cache
design preserves instruction selectivity that the end-to-end variant loses,
and that this difference is attributable to the trainable language adapter.

\section{Related Work}
\label{sec:related}

\subsection{Lightweight VLA policies}
Scale ranges from 7B-parameter OpenVLA-family models to sub-1B lightweights:
SmolVLA (0.45B) trains a compact flow head on VLM embeddings
\cite{shukor2025smolvlavisionlanguageactionmodelaffordable}; Evo-1 (0.77B)
shows two-stage training---frozen-VLM action-head pretraining followed by full
fine-tuning---consistently beats one-stage training \cite{lin2025evo1lightweightvisionlanguageactionmodel}; TurboVLA
(0.2B) removes the LLM trunk entirely and reaches 97.7\% LIBERO average
\cite{xie2026turbovlarealtimevisionlanguageactionmodel}. $\pi_0$ and $\pi_{0.5}$ (3B+) popularized flow-matching
action heads with pretrained VLM backbones \cite{black2026pi0visionlanguageactionflowmodel,intelligence2025pi05visionlanguageactionmodelopenworld}. Our work sits in
this lightweight family (43.5M trainable parameters, frozen Qwen-3.5-2B +
V-JEPA 2.1 extractors) and shares the flow-matching action head of
$\pi_0$/Evo-1/SmolVLA.

\subsection{Memory-augmented VLA}
ReMem-VLA maintains double-layer EMA recurrent queries with an
image-reconstruction auxiliary loss \cite{li2026rememvlaempoweringvisionlanguageactionmodel};
MemoryVLA builds a perceptual--cognitive memory bank with retrieval
\cite{shi2026memoryvlaperceptualcognitivememoryvisionlanguageaction};
AVA-VLA learns a recurrent belief with active visual attention under T=4 BPTT
\cite{xiao2026avavlaimprovingvisionlanguageactionmodels}; RB-VLA couples a
belief state with a world-model self-supervised objective
\cite{bagaria2026recursivebeliefvisionlanguage}; CogVLA couples visual and
action streams in a bidirectional attention decoder
\cite{li2026cogvlacognitionalignedvisionlanguageactionmodel}. Our VA
composite is closest to the recurrent bidirectional designs (ReMem/RB/AVA)
but differs in two respects: the memory is a single-slot overwritten visual
state per layer (constant memory, no retrieval or bank), and---central to this
paper---language enters through a \emph{frozen static cache} rather than a
trainable shared backbone, so our objective needs none of the auxiliary
losses these baselines use to protect representations.

\subsection{Instruction blindness and language grounding}
LIBERO provides multi-task same-scene manipulation with diverse instructions
\cite{liu2023liberobenchmarkingknowledgetransfer}; LIBERO-Plus formalizes robustness protocols including blank
and swapped instructions \cite{fei2025liberoplusindepthrobustnessanalysis};
LIBERO-CF introduces counterfactual benchmarks showing SOTA VLAs frequently
ignore language when scene features suffice
\cite{libero_cf}. On the training side, flatness/sharpness analyses
(SAM-based fine-tuning) report +217\% instruction-following on LIBERO-CF with
zero architecture change \cite{zhang2026flatnesspreservesinstructionfollowing}, and Knowledge Insulation
shows continuous action-expert gradients damage language-command
interpretation, motivating stop-gradient isolation of the backbone
\cite{driess2025knowledgeinsulatingvisionlanguageactionmodels}. Our frozen-cache design is the architectural
extreme of the stop-gradient direction: no gradient can reach language
parameters, verified directly (Sec.~\ref{sec:collapse}).

\subsection{Flow-matching policies and RL fine-tuning}
Flow matching and diffusion action heads are the de-facto standard
\cite{black2026pi0visionlanguageactionflowmodel,intelligence2025pi05visionlanguageactionmodelopenworld,nvidia2025gr00tn1openfoundation}. RL fine-tuning of such policies is addressed by
Diffusion Policy Policy Optimization (DPPO) \cite{ren2024diffusionpolicypolicyoptimization}, ReinFlow (policy
gradients through the denoising path) \cite{zhang2026reinflowfinetuningflowmatching} and $\pi$RL
(Flow-Noise/Flow-SDE augmented Markov policies; LIBERO few-shot SFT
57.6\%$\rightarrow$97.6\%, MetaWorld 50.8\%$\rightarrow$78.1\%/85.8\% with
sparse 0/1 success reward and 64 parallel environments)
\cite{pi_rl}. We follow the
ReinFlow-lite recipe: 32 learnable per-step per-dimension transition-noise
parameters on the imitation-learned flow head, PPO with the exact joint path
log-probability (Sec.~\ref{sec:rl}).

\section{Method}
\label{sec:method}

\begin{figure}[t]
\centering
\fbox{\parbox{0.9\columnwidth}{\centering\vspace{3em}
\textbf{Figure 1 (placeholder):} frozen Qwen cache $\rightarrow$ VA composite
(V/M/A streams, shared attention) $\rightarrow$ flow head.\vspace{3em}}}
\caption{Architecture overview. To be replaced by the final diagram.}
\label{fig:arch}
\end{figure}

\subsection{Architecture}
\label{sec:arch}
The policy has four components (Fig.~\ref{fig:arch}).
\begin{enumerate}
\item \textbf{Frozen language encoder with static cache.} The instruction
$l$ is encoded once by a frozen Qwen-3.5-2B text backbone. Final-layer hidden
states of all language tokens are cached as language memory $K_L, U_L$. At
deployment the encoder never runs per decision step---only when the
instruction changes. Critically, the encoder is frozen \emph{and} detached:
action gradients cannot flow into language parameters (Sec.~\ref{sec:mechanism}).
\item \textbf{Vision encoder.} V-JEPA 2.1 (ViT-B) encodes the current causal
4-frame window into visual tokens $V_t$; only the same-layer output of the
previous step is kept as target visual memory $M_{t-1}$, so history length
never grows.
\item \textbf{VA composite.} A stack of $N$ shared-attention blocks with
vision, action and memory streams plus static language anchors:
\begin{equation}
Q=\begin{bmatrix}V_tW_Q^V\\A_tW_Q^A\end{bmatrix},\quad
K=\begin{bmatrix}V_tW_K^V\\M_{t-1}W_K^M\\A_tW_K^A\\K_L\end{bmatrix},\quad
U=\begin{bmatrix}V_tW_U^V\\M_{t-1}W_U^M\\A_tW_U^A\\U_L\end{bmatrix},
\end{equation}
\begin{equation}
[V_t',A_t']=[V_t,A_t]+\operatorname{softmax}\left(\frac{QK^\top}{\sqrt d}\right)U .
\end{equation}
Each layer overwrites memory with $M_t=V_t'$. Vision queries read the
previous target visual memory, the action stream and the fixed language
anchor; action queries read the same. Memory is cleared when the command
changes. After $N$ layers the action stream is normalized into the semantic
action condition $C_t=\operatorname{LN}(A_t^N)$.
\item \textbf{Flow head.} A lightweight flow-matching module parameterizes
the vector field $\dot a^\tau=v_\theta(a^\tau,\tau\mid C_t)$ over the full
action chunk. Deployment runs the VA composite once; only the flow head
repeats across 32-step Euler integration.
\end{enumerate}

\subsection{Training}
\label{sec:training}
Backbones stay frozen during initial training; the VA composite is unrolled
over 4 consecutive visual timesteps so the action loss back-propagates through
$M_{t-1}$ to previous visual states, language projections, memory K/V and the
flow head. With $a^\tau=(1-\tau)\epsilon+\tau a$ and $\tau\sim U(0,1)$:
\begin{equation}
\LFM=\mathbb E\left\|v_\theta(a^\tau,\tau\mid C_t)-(a-\epsilon)\right\|_2^2 .
\end{equation}
Each training pair contains two mutually exclusive instructions at the same
0-th-step vision and robot state; the paired branch uses the same noise
$\epsilon_p$ at fixed $\tau=0$ so only the language condition differs:
\begin{equation}
\Lpair=\operatorname{Huber}\left[
v_\theta(\epsilon_p,0\mid C_i)-v_\theta(\epsilon_p,0\mid C_j),\;
a_i-a_j\right].
\end{equation}
Total loss $\mathcal L=\LFM+\lambda_{pair}\Lpair$.
\todo{Note: LIBERO e2e runs are FM-only (pair=0), itself a matched control.}

\subsection{Deployment}
Single VA forward + 32-step Euler integration of the flow head; 40.6~Hz on a
consumer GPU; no autoregressive decoding, no diffusion iteration.

\subsection{Why the frozen cache prevents representation erosion}
\label{sec:mechanism}
The causal chain behind instruction blindness is
\texttt{action gradients $\rightarrow$ shared backbone $\rightarrow$
language parameters rewritten}. In the VA composite this chain is cut at its
first link: the language encoder is called in \texttt{no\_grad} mode and its
outputs are detached before entering the VA composite, so the computation
graph contains no path into language parameters. This is a property of the
graph structure, not a training trick. The residual risk is not erosion of
the encoder but \emph{collapse of the read-out}---the VA's own language
projections could degenerate. We measure the encoder directly (pairwise
cosine of instruction embeddings, Sec.~\ref{sec:collapse}) and the read-out
behaviorally (counterfactual perturbations), keeping the two failure modes
distinguishable.

\section{Experimental Setup}
\label{sec:setup}
\textbf{Benchmarks.} (i)~PNPW single-task pick-and-place (118k frames) for
architecture ablations; (ii)~LIBERO 3-scene 12-task subset (360 samples),
the same-scene different-instruction regime that exposes visual shortcuts;
(iii)~MetaWorld MT50 (49 tasks $\times$ 50 demos, SmolVLA's
\texttt{lerobot/metaworld\_mt50} data source); \todo{LIBERO-100, CALVIN}.
Multi-start sampling: each MetaWorld demo contributes four training windows
at evenly spaced starts $\{0, L/3, 2L/3, L\}$ (where $L$ is the window
span), covering the full episode rather than only its first seconds.
\textbf{Protocol.} 32-step flow sampling; chunk MAE on normalized actions;
first-step success ($<0.05$); macro-average over tasks; 95\% bootstrap CI
(B=2000, seed 0, resampling unit = task). The copy-previous-action
persistence baseline is reported alongside every number
\cite{copycat_ptp}. \textbf{Language perturbations.} \emph{blank} (language
stream zeroed), \emph{swap} (instruction rotation +1), \emph{wrong}
(MetaWorld rotation +1 mod 49), \emph{task-id} (language replaced by a
task-identity token). Counterfactual closed-loop verdicts (L$_m$) use
same-scene dual-objective pairs where only the language cache differs
(Sec.~\ref{sec:obedience}).

\section{Results}
\label{sec:results}

\subsection{Single-task pilot (PNPW)}
\label{sec:pnpw}
Table~\ref{tab:pnpw}: depth must be matched with step budget: 8 layers
@20k beats 4 layers by 31\% chunk error (4$\times$ the real gain of the
4-layer config), and flat pooling (token-mean) beats spatial pooling. An
intervention audit showed the action$\rightarrow$vision backward path
(+7.7\% error), memory$\rightarrow$action (+2.9\%) and previous-action
(dominant input) are all real; language showed $\sim$0\% because the
instruction is constant in single-task data.

\begin{table}[t]
\centering\small
\caption{PNPW ablations.}
\label{tab:pnpw}
\begin{tabular}{lcccc}
\toprule
config & success & first-mae & chunk & vs.\ persist. \\
\midrule
4L @10k (flat) & 90.9\% & 1.228 & 0.0497 & $-0.0055$\\
4L @10k (spatial) & 84.5\% & 1.542 & 0.0564 & $-0.0014$\\
8L @10k (flat) & 85.7\% & 1.406 & 0.0518 & $-0.0034$\\
8L @20k (flat) & 99.1\% & 1.063 & 0.0345 & $-0.0207$\\
persistence & 99.7\% & --- & 0.0552 & ---\\
\bottomrule
\end{tabular}
\end{table}

\subsection{Language grounding on LIBERO (open loop)}
\label{sec:libero}

\textbf{Reading.} In the same-scene different-instruction regime---where the
scene alone cannot disambiguate the goal---blanking the language stream
inflates chunk error by 2381\% (3-scene) and swapping instructions by 607\%,
confirming language is a \emph{necessary} input, not a weak conditioner. The
1-scene numbers are even sharper (+13751\% / +1518\%) because four
instructions share one scene and the policy has no other disambiguating
signal; the 3-scene set includes scenes where multiple tasks differ
(drawer-open vs.\ bowl placement) and the visual state carries more task
identity, dampening the perturbation effect. This monotonicity---sensitivity
rises as the visual shortcut signal weakens---is exactly the signature of
genuine language grounding rather than spurious correlation.

Table~\ref{tab:trio}: language is causally necessary under same-scene
different-instruction data---the regime where VLA shortcut risk is highest.

\begin{table}[t]
\centering\small
\caption{LIBERO language trio (32-step).}
\label{tab:trio}
\begin{tabular}{lccc}
\toprule
dataset & clean chunk & blank & swap \\
\midrule
1-scene (4 tasks) & 0.00106 & +13751\% & +1518\%\\
3-scene (12 tasks) & 0.00254 & +2381\% & +607\%\\
\bottomrule
\end{tabular}
\end{table}

\subsection{MetaWorld MT50}
\label{sec:mw}
Open loop (multi-start full-coverage data, 49$\times$4 decision points):
chunk MAE 0.0806 [0.0709, 0.0921] vs.\ persistence 0.0930 ($-14\%$);
first-step success 50.8\% [45.3\%, 56.1\%] (persistence-dominated,
reported per Sec.~\ref{sec:setup}).

\textbf{Language vs.\ task identity.} Rotating the instruction (+1 mod 49)
inflates chunk error by +108.5\%; replacing language with a task-identity token
("task $i$") inflates it by +81.5\%. Two readings: the task-id gap is large
because MetaWorld tasks are partly identifiable from the visual state, but the
\emph{extra} degradation of wrong-instruction over task-id (+108.5\% vs.\
+81.5\%) is attributable to language content beyond task identity---the
policy uses instruction semantics, not merely a task label. This is the
counterpart, on MT50, of the LIBERO verdict (Sec.~\ref{sec:libero}):
language is grounded when it is the only disambiguating signal, and still
causally relevant when it is not.

Closed loop (49$\times$10, horizon 500): multi-start rebuild raises success
from 7.1\% to \textbf{16.3\%} [9.4\%, 24.1\%] (80/490). The remaining gap to
closed-loop literature numbers (SmolVLA 0.45B 57.3\%, Evo-1 80.6\%,
10-trials protocol) is dominated by two factors we verify rather than assert:
(i) the trained policy's coverage is still partial---four phase-local windows
per demo---so long-horizon tasks (nut-peg, stick-push, sweep, soccer) roll
out almost entirely out of distribution (0--3/10), while short tasks succeed
near ceiling (button 8/10, door-close 8/10, drawer-close 8/10,
window-close 9/10); and (ii) closed-loop action errors accumulate through
the previous-action/visual-memory recurrence (Sec.~\ref{sec:mechanism}).
Sec.~\ref{sec:rl} reports the VLA-RL fine-tuning that closes part of this
gap; per-task numbers are in the supplementary logs.

\subsection{Comparison to literature baselines}
\label{sec:baselines}
Table~\ref{tab:baselines}: all baseline numbers cited from original papers
(Grok-verified protocols); no baselines trained by us.

\begin{table*}[t]
\centering\small
\caption{Literature comparison. All baseline numbers cited from the original
papers or the verified downstream reports listed in the footnotes; no baselines
trained by us. Protocols noted per row family.}
\label{tab:baselines}
\begin{tabular}{lcccl}
\toprule
model & MetaWorld MT50 (closed-loop) & LIBERO avg & params & notes \\
\midrule
FabriVLA & 90.0\% & --- & --- & imitation SOTA\\
LA4VLA & 87.5\% & --- & --- & \\
$\pi_0$+ALAM & 85.0\% & --- & --- & \\
Evo-Depth & 84.4\% & --- & --- & \\
Evo-1 & 80.6\%$^{\dagger}$ & 94.8\%$^{\dagger}$ & 0.77B & two-stage; 10 trials$\times$5 seeds\\
SmolVLA 0.45B & 57.3\%$^{\dagger}$ & 87.3\%$^{\dagger}$ & 0.45B & 10 trials/task\\
SmolVLA 2.25B & 68.24\%$^{\dagger}$ & 88.75\%$^{\dagger}$ & 2.25B & \\
TurboVLA & --- & 97.7\%$^{\ddagger}$ & 0.2B & no LLM trunk\\
$\pi_{0.5}$ & --- & 96.9\%$^{\S}$ & 3B+ & OpenPI full-SFT report\\
$\pi_0$ & 47.9\%$^{\P}$ & 94.2\%$^{\S}$ & 3B+ & both third-party\\
OpenVLA & --- & 76.5\%$^{\text{\S\P}}$ & 7B & Appendix E (v2)\\
$\pi$RL ($\pi_0$ SFT$\rightarrow$Flow-Noise) & 50.8$\rightarrow$85.8\%$^{\|}$ & --- & 3B+ & sparse 0/1 reward, 64 envs\\
\midrule
ours (open-loop) & chunk 0.0806 & --- & 43.5M & not comparable\\
ours (closed-loop) & 13.9\% [7.6, 21.0] & --- & 43.5M & macro4 = 7.0\%; gap analysis \S\ref{sec:mw}\\
\bottomrule
\end{tabular}

\smallskip
\footnotesize
$^{\dagger}$ MT50: 50 demos/task, 10 trials/task (Seo split); VLM-init only.
Evo-1 averages 5 seeds; SmolVLA per Table 2.
$^{\ddagger}$ LIBERO: TurboVLA Table 1, 50 rollouts/task (2000 trials),
VLA-Adapter rollout protocol.
$^{\S}$ $\pi_{0.5}$/$\pi_0$ LIBERO numbers are \emph{not} from the original
papers (which report no LIBERO experiments); they are OpenPI full-dataset SFT
results as tabulated in $\pi$RL (Table 1) / TurboVLA (Table 1), 50 trials/task
(10 subtasks $\times$ 50 states = 500 rollouts/suite).
$^{\P}$ $\pi_0$ MetaWorld 47.9\% is a third-party citation (SmolVLA Table 2).
$^{\text{\S\P}}$ OpenVLA LIBERO is the authors' Appendix E (v2) fine-tuning
report, 3 seeds $\times$ 500 rollouts/suite (10$\times$50); the main paper
reports no LIBERO experiments.
$^{\|}$ $\pi$RL MetaWorld results per the RLinf/$\pi$RL MetaWorld eval
protocol (SFT on \texttt{lerobot/metaworld\_mt50}, 2500 trajectories, sparse
0/1 success reward, 64 parallel envs); SmolVLA 68.2\% therein is a
third-party citation.
\end{table*}

\subsection{End-to-end fine-tuning collapse and the 2$\times$2 control}
\label{sec:collapse}

\textbf{Reading of the collapse.} The cosine verdicts decouple two hypotheses
about why B40k loses language sensitivity. If fine-tuning with a trainable
language adapter (\emph{any} adapter) were benign, the instruction embedding
space would stay close to the pretrained one while the read-out changed;
instead the space itself collapses (0.999 vs.\ 0.765). If the collapse were
driven by vision fine-tuning alone, C2 (vision frozen, LoRA only) would keep
the pretrained space; if driven by the trainable language adapter, C2 would
collapse like B40k. The combination of the representation-level verdict
(cosine) and the behavior-level verdicts (blank/swap sensitivity, C\_OL
displacement) therefore pins the driver of instruction blindness within our
settings to the trainable adapter, and motivates the repair directions we
evaluate: language-side isolation ($\Lpair$ reweighting / frozen adapter) or
visual anchoring (regularization toward the frozen encoder's features).
C2 (vision frozen, LoRA only) reproduces the collapse (cosine 0.999), confirming that the trainable language adapter---not visual fine-tuning---drives the collapse within our settings; language-side isolation and visual anchoring remain the repair directions.
Table~\ref{tab:2x2}: B40k (V-JEPA 12 blocks + Qwen LoRA r32) collapses the
instruction embedding space to near-identity (cosine 0.999) and loses
behavioral sensitivity (+1.5\% vs.\ +2381\%). C1---the same video-input e2e
pipeline with all backbones frozen---keeps the exact original embedding space
(0.765), isolating the collapse to \emph{training} rather than the e2e
pipeline itself.

\begin{table}[t]
\centering\small
\caption{2$\times$2 control: embedding cosine and behavioral sensitivity.}
\label{tab:2x2}
\begin{tabular}{lcccc}
\toprule
config & V-JEPA & Qwen & cosine & blank sens.\\
\midrule
A (frozen) & frozen & frozen & 0.765 & +2381\%\\
B40k (e2e) & 12 blk & LoRA r32 & 0.999 & +1.5\%\\
C1 (e2e frozen) & frozen & frozen & 0.765 & +33.0\%\\
C2 (LoRA only) & frozen & LoRA r32 & 0.999 & \todo{}\\
random feats & --- & --- & 0.083 & ---\\
\bottomrule
\end{tabular}
\end{table}

\subsection{Language obedience: counterfactual closed-loop verdict}
\label{sec:obedience}
Same-scene dual-objective rollouts: on each matched initial state we roll out
four conditions where only the language cache differs:
$D$: env($g_1$)+$l_1$, env($g_2$)+$l_2$; $O$: env($g_1$)+$l_2$,
env($g_2$)+$l_1$. $L_m = D-O$ with block bootstrap 95\% CI. Verdicts:
$L_m\gg0$ obedience; $L_m\approx0$ with high $D,O$ missing language
selectivity; $L_m\approx0$ with low $D,O$ OOD fragility. Pairs: STUDY
back/front and left/right caddy compartments; KITCHEN black bowl
back/front; LIVING soup/butter and milk/juice (5 pairs, benchmark-pinned).

\textbf{Why matched blocks.} Because $D$ and $O$ conditions share the same
physical environment, initial state and sampling noise (only the language
cache differs), any difference in success is attributable to language---this
is a within-episode counterfactual, not a between-task comparison. The
swapped conditions double as a held-out command-fork test: the same scene is
shown with a different but executable instruction, exposing visual shortcuts
that a clean-task metric would miss. \todo{numbers for A and B40k}

\subsection{VLA-RL}
\label{sec:rl}
Following $\pi$RL/ReinFlow, we fine-tune the IL policy with PPO over a sparse
success reward, freezing V-JEPA and Qwen and training only the VA composite,
flow head, a value head and a 32-parameter flow-noise schedule. The policy is
the augmented Markov flow policy: each Euler transition draws
$x_{k+1}\sim\mathcal N(x_k+\frac1K v_\theta(\cdot),\,\sigma_k^2 I)$ with
$\sigma_{k,d}=0.02+0.06\,\mathrm{sigmoid}(\alpha_{k,d})$. PPO uses the exact
joint path log-probability; the transition noise is dropped at evaluation.
Macro actions: 8-step chunks executed for the first 6 primitives; sparse
success reward; GAE ($\gamma=0.99,\lambda=0.95$), 4 epochs, minibatch 128,
clip 0.1, actor LR $3\times10^{-6}$, critic LR $10^{-4}$.
\todo{MT10 IL$\rightarrow$RL numbers}

\subsection{Efficiency}
\label{sec:efficiency}
Trainable parameters 43.5M (4-layer VA + flow head); deployment latency
24.66~ms/decision (40.6~Hz at 32 Euler steps); training on a single 24GB GPU;
baselines: SmolVLA 0.45B total, $\pi_0$ 3B+ (dual A100).

\section{Real-Robot Validation}
\label{sec:real}
\textbf{(Structure reserved; experiments to be conducted by the authors on
hardware.)} \todo{To be filled: robot platform, camera setup, task suite,
success criteria, and sim-to-real observations.} The frozen-cache design has
two properties relevant to deployment: the language encoder runs once per
instruction (no per-step LLM inference), and the VA composite + flow head
deploys at 40.6~Hz on embedded hardware \todo{measured on the target
platform}. We plan to validate: (i)~instruction following under swapped and
novel instructions; (ii)~closed-loop robustness over extended horizons; and
(iii)~the counterfactual verdicts of Sec.~\ref{sec:obedience} on physical
rollouts.

\section{Limitations}
\label{sec:limitations}
\begin{itemize}
\item \textbf{Closed-loop data coverage.} Our MetaWorld training samples
covered only the first 0.33~s of each episode, so closed-loop rollouts ran
mostly out of distribution (7.1\% [2.7\%, 12.7\%]); in-distribution tasks
succeed at ceiling, attributing the gap to data coverage. The multi-start
rebuild (4 segments per episode) is designed to close this gap
\todo{results}. We do not claim parity with closed-loop literature numbers
until that retest.
\item \textbf{LIBERO closed-loop alignment.} The kevin\_libero100 data lacks
object-pose fields, so official-environment init alignment is not possible
from the data; LIBERO closed-loop evidence relies on the same-scene
counterfactual protocol (Sec.~\ref{sec:obedience}) instead of the standard
benchmark rollouts.
\item \textbf{Protocol mixing in comparisons.} Literature numbers are cited
with per-family footnotes (10 vs.\ 50 trials); our open-loop numbers are not
directly comparable to closed-loop baselines and are labeled as such.
\item \textbf{Additional benchmarks.} CALVIN data was unreachable from our
network (dataset server); RLBench requires CoppeliaSim/PyRep (a week-scale
port) and is deferred. RoboTwin 2.0 is bimanual with 14-D actions, not
directly supported by the single-arm architecture.
\item \textbf{Real-robot validation} is planned (Sec.~\ref{sec:real}) and
will be reported separately.
\end{itemize}

\section{Conclusion}
\label{sec:conclusion}
The VA composite removes representation erosion by construction: action
gradients never touch language parameters. Our $2\times2$ control confirms
both halves---freezing preserves the pretrained instruction embedding space
exactly (cosine 0.765) and behavioral language sensitivity (+2381\%), while
end-to-end LoRA fine-tuning collapses the space (0.999) and removes
selectivity (+1.5\%). Counterfactual evaluation on MetaWorld confirms
language content beyond task identity is used (wrong-instruction +108.5\%).
We remain scoped: visual shortcuts are a property of the data distribution
that no architecture alone removes, and our closed-loop numbers are currently
limited by training-data coverage, not architecture.
\todo{final numbers: multi-start, L\_m, VLA-RL, LIBERO-100}

\bibliographystyle{abbrv}
\bibliography{references}

\end{document}
